\documentclass[10pt]{article}
\providecommand{\FIGDIR}{fig}
\newcommand{\DIGESTDATE}{30 July 2026}
% ===================== PM Research Watch - shared preamble =====================
\usepackage[T1]{fontenc}
\usepackage[utf8]{inputenc}
\usepackage{amsmath,amssymb}
\usepackage{tgpagella}
\usepackage{tgheros}
\renewcommand{\sfdefault}{qhv}
\renewcommand{\ttdefault}{lmtt}
\usepackage{microtype}
\usepackage[english]{babel}

\usepackage[a4paper,top=17mm,bottom=17mm,left=17mm,right=17mm,
            headsep=6mm,footskip=10mm]{geometry}
\usepackage{graphicx}
\usepackage{xcolor}
\usepackage{colortbl}
\usepackage{tikz}
\usetikzlibrary{calc}
\usepackage[most]{tcolorbox}
\usepackage{titlesec}
\usepackage{fancyhdr}
\usepackage{booktabs}
\usepackage{array}
\usepackage{longtable}
\usepackage{enumitem}
\usepackage{multicol}
\usepackage{ragged2e}
\usepackage[hidelinks]{hyperref}

% ------------------------------------------------------------------ palette
\definecolor{Ink}    {HTML}{14262E}
\definecolor{Deep}   {HTML}{0F4C5C}
\definecolor{Teal}   {HTML}{1F7A8C}
\definecolor{Sky}    {HTML}{4FA3C4}
\definecolor{Sage}   {HTML}{6FA287}
\definecolor{Amber}  {HTML}{E3A008}
\definecolor{Clay}   {HTML}{B5643C}
\definecolor{Coral}  {HTML}{D1495B}
\definecolor{Violet} {HTML}{6A5B9E}
\definecolor{Slate}  {HTML}{4A6572}
\definecolor{Paper}  {HTML}{FBF9F5}
\definecolor{Mist}   {HTML}{EEEAE2}
\definecolor{Rule}   {HTML}{D9D3C9}

\pagecolor{Paper}
\color{Ink}

\hypersetup{colorlinks=true, linkcolor=Deep, urlcolor=Teal, citecolor=Teal,
            pdftitle={PM Research Watch}, pdfauthor={Automated daily digest}}

% ------------------------------------------------------------------ headings
\titleformat{\section}
  {\sffamily\bfseries\large\color{Deep}}{}{0pt}
  {}[{\vspace{1.2mm}\color{Amber}\titlerule[1.1pt]}]
\titlespacing*{\section}{0pt}{5.5mm}{2.8mm}

\titleformat{\subsection}
  {\sffamily\bfseries\normalsize\color{Teal}}{}{0pt}{}
\titlespacing*{\subsection}{0pt}{3.6mm}{1.4mm}

\setlist[itemize]{leftmargin=4.2mm, itemsep=1.1mm, topsep=1.3mm, parsep=0pt,
                  label=\textcolor{Teal}{\footnotesize$\blacktriangleright$}}

% ------------------------------------------------------------------ header/footer
\pagestyle{fancy}
\fancyhf{}
\renewcommand{\headrulewidth}{0pt}
\renewcommand{\footrulewidth}{0pt}
\fancyhead[L]{\sffamily\scriptsize\color{Slate}PM RESEARCH WATCH\;\textcolor{Rule}{\textbar}\;\DIGESTDATE}
\fancyhead[R]{\sffamily\scriptsize\color{Slate}Particulate matter monitoring \& health effects}
\fancyfoot[C]{\sffamily\scriptsize\color{Slate}\thepage}

% ------------------------------------------------------------------ components
% metric tile
\newcommand{\metric}[3]{%
  \begin{tcolorbox}[enhanced, colback=white, colframe=Rule, boxrule=0.5pt,
    arc=1.6mm, left=2.2mm, right=2.2mm, top=1.8mm, bottom=1.8mm,
    borderline west={1.7mm}{0pt}{#3}, nobeforeafter, valign=center,
    height=17mm, width=\linewidth]
    \RaggedRight
    {\sffamily\bfseries\fontsize{17}{17}\selectfont\color{#3}#1}\\[0.5mm]
    {\sffamily\fontsize{6.5}{7.4}\selectfont\color{Slate}#2}
  \end{tcolorbox}}

% pull-quote / key insight box
\newtcolorbox{keybox}[1][]{enhanced, colback=Mist, colframe=Mist,
  arc=1.6mm, left=3.2mm, right=3.2mm, top=2.4mm, bottom=2.4mm,
  borderline west={1.7mm}{0pt}{Amber}, #1}

% subtopic band
\newcommand{\band}[2]{%
  \vspace{3.2mm}\noindent
  \begin{tcolorbox}[enhanced, colback=#2, colframe=#2, arc=1.2mm,
    left=2.6mm, right=2.6mm, top=1.1mm, bottom=1.1mm, boxrule=0pt,
    width=\linewidth, nobeforeafter]
    {\sffamily\bfseries\small\color{white}#1}
  \end{tcolorbox}\vspace{1.6mm}}

% one paper entry: {short cite}{title}{journal}{doi}{body}
\newcommand{\paperentry}[5]{%
  \par\vspace{1.9mm}\noindent
  {\sffamily\bfseries\footnotesize\color{Ink}#1}\;%
  {\sffamily\fontsize{7}{8}\selectfont\color{Slate}#3}\\[0.4mm]
  {\footnotesize\itshape\color{Deep}#2}\\[0.6mm]
  {\footnotesize #4}\\[0.3mm]
  {\fontsize{6.6}{7.6}\selectfont\color{Slate}\texttt{doi:}\,\href{https://doi.org/#5}{\texttt{\detokenize{#5}}}}
  \par}

% caption style
\newcommand{\figcap}[1]{\par\vspace{1.0mm}
  {\fontsize{7.2}{8.6}\selectfont\color{Slate}\textsf{\textbf{Fig.}}\; #1\par}\vspace{1.5mm}}

\providecommand{\permil}{\textperthousand}

\setlength{\parindent}{0pt}
\setlength{\parskip}{0pt}
\renewcommand{\arraystretch}{1.18}


\begin{document}

% ==================================================================== MASTHEAD
\thispagestyle{empty}
\begin{tikzpicture}[remember picture, overlay]
  \fill[Deep] (current page.north west) rectangle
      ($(current page.north east)+(0,-67mm)$);
  \fill[Amber] ($(current page.north west)+(0,-67mm)$) rectangle
      ($(current page.north east)+(0,-68.4mm)$);
  \foreach \x/\y/\r in {22/12/0.9, 41/26/1.5, 63/9/0.7, 88/31/1.1, 112/17/1.9,
                        138/29/0.8, 159/13/1.3, 178/54/1.0, 196/21/1.6,
                        30/44/0.6, 74/57/1.2, 125/48/0.9, 168/6/0.8, 191/59/0.7,
                        52/18/1.0, 100/60/0.6, 148/20/1.1, 66/49/1.4, 205/45/0.8}
    \fill[Sky, opacity=0.30] ($(current page.north west)+(\x mm,-\y mm)$) circle (\r mm);
\end{tikzpicture}

\vspace*{4mm}
{\sffamily\fontsize{7.6}{9}\selectfont\color{Sky}%
 AUTOMATED DAILY LITERATURE DIGEST\;\textbullet\;PhD RESEARCH WATCH\par}
\vspace{2.2mm}
{\sffamily\bfseries\fontsize{27}{29}\selectfont\color{white}%
 Particulate Matter\par}
{\sffamily\fontsize{27}{29}\selectfont\color{Sky}%
 Monitoring \& Health Effects\par}
\vspace{5.0mm}
{\sffamily\fontsize{8.4}{10}\selectfont\color{white}%
 Issue 2026-07-30\;\;\textcolor{Amber}{\textbar}\;\;Entry window 30 July 2026\;\;%
 \textcolor{Amber}{\textbar}\;\;18 records in scope\par}
\vspace{18mm}

% -------------------------------------------------------------- metric strip
\noindent
\begin{minipage}[t]{0.190\linewidth}\metric{18}{RECORDS IN SCOPE\\(9 SUBTOPICS)}{Deep}\end{minipage}\hfill
\begin{minipage}[t]{0.190\linewidth}\metric{10}{EFFECT ESTIMATES\\EXTRACTED}{Teal}\end{minipage}\hfill
\begin{minipage}[t]{0.190\linewidth}\metric{6}{RECORDS FROM CHINA\\(33\% OF CORPUS)}{Sky}\end{minipage}\hfill
\begin{minipage}[t]{0.190\linewidth}\metric{4}{SOUTH ASIA RECORDS\\(2 PREPRINTS)}{Amber}\end{minipage}\hfill
\begin{minipage}[t]{0.190\linewidth}\metric{1}{SENSING /\\INSTRUMENTATION}{Clay}\end{minipage}

\vspace{6mm}

% -------------------------------------------------------------- signal of day
\section*{\sffamily\bfseries\color{Deep}Signal of the day}
\vspace{-1.5mm}{\color{Amber}\rule{\linewidth}{1.1pt}}\vspace{2.5mm}

\begin{keybox}
{\sffamily\bfseries\small\color{Deep}Aerosol liquid water content is now available as a
1\,km daily gridded field for the contiguous US --- the missing thermodynamic covariate
in both sensor calibration and toxicity-weighted exposure.}\\[1.4mm]
\footnotesize
Zhang et al.\ (\textit{npj Clim Atmos Sci}) trained machine-learning surrogates on
GEOS-Chem output to reproduce the thermodynamic relationship between ALWC and its
predictors, then applied those relationships to bias-corrected high-resolution inputs,
yielding a \textbf{1\,km $\times$ 1\,km daily ALWC field for 2000--2019}. The ML product
resolves daily variation and spatial heterogeneity that the parent CTM does not; ALWC is
highest in the Midwest and lowest in the West, and \textbf{declined across most regions
over the period, driven primarily by sulfate reduction}. The authors then show ALWC
constrains variability in water-soluble iron, a health-relevant redox-active metal
fraction.

For distributed low-cost PM sensing this matters twice over. First, ALWC is the physical
quantity that hygroscopic-growth corrections are proxying for when they regress on
relative humidity; a gridded ALWC field is a strictly better covariate than RH alone,
because it already carries the composition dependence that makes RH corrections
non-transferable between regions. Second, the sulfate-driven secular decline means a
correction function fitted on a 2015 co-location is being applied to an aerosol whose
water content has since shifted --- a mechanism for calibration drift that is
compositional rather than electronic, and which no amount of sensor-side quality control
will detect.
\end{keybox}

\vspace{3mm}

% -------------------------------------------------------------- exec brief
\section*{\sffamily\bfseries\color{Deep}What changed today}
\vspace{-1.5mm}{\color{Amber}\rule{\linewidth}{1.1pt}}\vspace{2.5mm}

\footnotesize
\begin{itemize}
\item \textbf{The instrumentation axis returned almost nothing this window --- and that
is a measurement about the pipeline, not the field.} PubMed does not index
\textit{Atmos.\ Meas.\ Tech.} or \textit{Atmos.\ Chem.\ Phys.}, and the Consensus sweep
that exists to close that gap has \emph{no entry-date filter}: every low-cost-sensor
paper it surfaced resolved to a January--July 2026 publication date, and one
(Ginsburg et al., \textit{Aerosol Sci Technol}) was already carried in the 29 July issue.
All eight were rejected rather than passed off as new. See \textit{Provenance}.

\item \textbf{Source apportionment reached the shipping sector with isotopic evidence.}
Chen et al.\ attribute \textbf{17.6--37.9\% of PM$_{2.5}$ nitrate} across five Chinese
coastal cities to shipping, using a dual-isotope Bayesian mixing model on 615 daily
samples --- with the striking regional inversion that shipping contribution
\emph{rises} with pollution severity in northern cities and \emph{falls} in southern
ones. Shipping NO$_x$ standards remain looser than road-vehicle standards; this is the
quantitative case against that asymmetry.

\item \textbf{A clean negative-direction effect estimate in the reproductive literature.}
Li et al.\ report full-cycle O$_3$ associated with a \textbf{15.45\% reduction in sperm
concentration (95\% CI $-$20.95 to $-$9.58)}, sub-staged across the meiotic window with
the peak at days 31--45. The design --- quarterly panel over a full 90-day spermatogenic
cycle in 18--22-year-olds --- is unusually well matched to the biology. The sample is
32 men; treat the window resolution as the contribution and the point estimates as
provisional.

\item \textbf{Two mid-size Chinese time-series studies landed with tight intervals.}
Nanning CVD mortality (87{,}913 deaths, 11 years) and Guangdong eczema outpatient visits
(879{,}356 records) both report small but precisely bounded per-unit associations, with
SO$_2$ carrying the largest single-pollutant coefficient in each. The eczema mixture
analysis attributes 49\% of the PM$_{10}$-based mixture effect to PM$_{10}$ itself, but
38\%/38\% of the PM$_{2.5}$-based mixture to SO$_2$ and NO$_2$ --- i.e.\ PM$_{2.5}$ mass
is close to a passenger variable in that model.

\item \textbf{South Asia is over-represented today, and by preprints.} Four of 18 records
are South Asian, two of them Research Square preprints reporting single-station
descriptive PM statistics (Cumilla, Bangladesh: February mean PM$_{10}$
210.7\,$\mu$g/m$^3$). Neither is peer reviewed; neither reports instrument provenance or
QA. Logged for coverage, not for citation.

\item \textbf{Conflict as an exposure event was quantified.} Abi Nader et al.\ paired
satellite and five-site ground PM$_{2.5}$ across the 2024 Beirut conflict; satellite
median rose 15.11 $\rightarrow$ 20.21\,$\mu$g/m$^3$, and BAPHE exposure--response
functions give RR 1.011 respiratory and 1.032 skin admissions. The transferable part is
that a ten-day handheld campaign can anchor a satellite comparison.
\textbf{Trial watch: zero} ClinicalTrials.gov registrations or updates matched the
window; the 23--29 July accumulation reports in the 2 August weekly rollup.
\end{itemize}

\clearpage

% ==================================================================== ANALYTICS
\section*{\sffamily\bfseries\color{Deep}Corpus analytics}
\vspace{-1.5mm}{\color{Amber}\rule{\linewidth}{1.1pt}}\vspace{3mm}

\noindent\includegraphics[width=\linewidth]{\FIGDIR/f1_subtopics.png}
\figcap{Subtopic assignment of the 18 in-scope records. Assignment is single-label by
dominant contribution. \textit{Exposure assessment \& modelling} and \textit{Burden,
policy \& mitigation} tie for the lead at four records each; \textit{Mechanistic
toxicology} is empty for the first time in this series, and \textit{Sensing,
forecasting \& instrumentation} carries a single record --- a commentary letter rather
than primary work.}

\vspace{2mm}
\noindent\includegraphics[width=\linewidth]{\FIGDIR/f2_design_endpoint.png}
\figcap{Left --- study architecture. Reviews and syntheses account for 4/18 (22\%) and
measurement campaigns for a further 4; primary observational epidemiology (cohort,
acute, cross-sectional) contributes 6. Right --- health endpoint by organ system.
Five records carry no human health endpoint at all --- the four exposure-assessment
records plus the forecasting commentary.}

\vspace{2mm}
\noindent\includegraphics[width=\linewidth]{\FIGDIR/f3_metric_geography.png}
\figcap{Left --- particle metric. PM$_{2.5}$ alone is the exposure in 7 records and
appears jointly with PM$_{10}$ in a further 4; five records treat PM as unspeciated mass
or as an emissions category, characteristic of the review and burden-estimate literature.
\textbf{No record today addresses ultrafine particles} --- the third consecutive issue in
which the deposition-relevant size fraction is unrepresented. Right --- geographic
provenance. China contributes 6 records (33\%), South Asia 4.}

\vspace{2mm}
\noindent\includegraphics[width=\linewidth]{\FIGDIR/f6_lifecourse.png}
\figcap{Life-course windows addressed. Counts are non-exclusive and cover only the seven
records with an age-specific human population; the remaining eleven are monitoring,
modelling, burden-estimate or review work with no defined life stage.
\textbf{No record this window addresses preconception or in-utero exposure} --- a break
from the developmental-origins emphasis that dominated the 27--29 July issues.}

\vspace{2mm}
\noindent\includegraphics[width=\linewidth]{\FIGDIR/f4_heatmap.png}
\figcap{Subtopic $\times$ study-architecture cross-tabulation. The structural feature this
series keeps flagging is starkest today: the \textit{Exposure assessment} and
\textit{Sensing} rows carry no health-endpoint linkage at all, while every
health-endpoint row is populated exclusively by designs that take exposure as a given
input.}

\vspace{2mm}
\noindent\includegraphics[width=\linewidth]{\FIGDIR/f5_forest.png}
\figcap{All quantitative effect estimates reported across the corpus, on a common
logarithmic axis. \textbf{Exposure contrasts are heterogeneous and estimates are not
mutually comparable}: Luo et al.\ report per-10\,$\mu$g/m$^3$ single-day-lag RRs;
Li et al.\ report per-IQR ratios, plotted below 1 because the association is a
\emph{reduction} in sperm concentration; Ouyang et al.\ mix per-10\,$\mu$g/m$^3$ and
per-quartile-of-mixture contrasts; Abi Nader et al.\ report a period contrast.
Whiskers are drawn only where a confidence interval was reported --- four of the ten
estimates today carry none, which is itself a reporting-quality signal.}

\clearpage

% ==================================================================== DIGEST
\section*{\sffamily\bfseries\color{Deep}Paper-level digest}
\vspace{-1.5mm}{\color{Amber}\rule{\linewidth}{1.1pt}}\vspace{1mm}

{\fontsize{7.4}{9}\selectfont\color{Slate}Ordered by cluster. Each entry gives the
methodological core and the finding that matters, not an abstract paraphrase.\par}

\band{Exposure assessment \& modelling\hfill 4 records}{Amber}

\paperentry{Zhang et al.\ 2026}
{High-resolution aerosol liquid water content across the contiguous US using machine learning}
{npj Clim Atmos Sci}
{ML surrogates trained on GEOS-Chem output to learn ALWC--predictor thermodynamics, then
applied to bias-corrected high-resolution inputs, giving \textbf{1\,km daily ALWC for
2000--2019}. Outperforms the parent CTM on daily variability and spatial heterogeneity.
Regional ordering Midwest $>$ East $>$ West is driven by PM$_{2.5}$ mass, composition,
temperature and RH; the secular decline tracks sulfate. Demonstrated as a constraint on
water-soluble iron variability. The limitation is structural: the ML model inherits every
bias of GEOS-Chem's thermodynamic module, so this is a downscaling of a simulation, not
an observational product --- there is no independent ALWC validation set. Still the most
directly useful record here for hygroscopic-growth correction of optical sensors.}
{10.1038/s41612-026-01371-2}

\paperentry{Chen et al.\ 2026}
{Dual-isotope evidence for shipping contributions to atmospheric nitrate in coastal Chinese cities}
{J Hazard Mater}
{615 daily PM$_{2.5}$ samples across Tianjin, Qingdao, Shanghai, Xiamen and Guangzhou
during the cold-season transition, with nitrate $\delta^{15}$N and $\delta^{18}$O fed to
a Bayesian isotope mixing model against an emissions-inventory prior. Shipping-derived
nitrate ranged \textbf{17.6$\pm$4.7\% (Tianjin) to 37.9$\pm$11.4\% (Xiamen)}. Northern
cities show shipping contribution rising with pollution severity and marine air-mass
association; southern cities show it falling as inland transport intensifies. The
inference threat is the mixing-model prior --- source endmember $\delta^{15}$N ranges for
shipping overlap substantially with those for coal combustion, and no endmember
sensitivity analysis is reported. Relevant to sensing because it quantifies how much of
coastal PM$_{2.5}$ mass a mass-only monitor cannot attribute.}
{10.1016/j.jhazmat.2026.143129}

\paperentry{Patwary 2026}
{Spatio-temporal assessment of PM$_{10}$ and meteorological correlates in Cumilla, Bangladesh}
{Res Sq (preprint)}
{Single Department of Environment monitoring station, January--April 2026, Pearson
correlation and trend analysis against temperature, RH and rainfall. Dry-season peak
\textbf{210.7\,$\mu$g/m$^3$ monthly mean PM$_{10}$ in February}, roughly 4.7$\times$ the
WHO 24-h guideline. No instrument model, calibration history or data-completeness
statistics are reported, and one station cannot support the ``spatio-temporal'' framing
in the title. Logged because South Asian monitoring records are structurally scarce in
this corpus, not because the analysis is sound.}
{10.21203/rs.3.rs-10519411/v1}

\paperentry{Raj \& Dhada 2026}
{Spatiotemporal ambient air quality across urban, tourist and industrial Himachal Pradesh}
{Res Sq (preprint)}
{Five Himalayan towns (Solan, Baddi, Shimla, Kullu, Manali) contrasted by land use,
25 June--24 July 2026, 150 site-days, using AQI plus PM$_{2.5}$, PM$_{10}$, CO,
SO$_2$, NO$_2$ and O$_3$ taken from the public AQI.in platform. The land-use contrast
(industrial Baddi vs.\ tourist hill towns) is the design's one genuine idea. The fatal
weakness is provenance: AQI.in aggregates heterogeneous and partly low-cost sources with
no published calibration chain, so between-town differences cannot be separated from
between-instrument differences. Preprint, not peer reviewed.}
{10.21203/rs.3.rs-10492149/v1}

\band{Burden, policy \& mitigation\hfill 4 records}{Coral}

\paperentry{Abi Nader et al.\ 2026}
{Short-term respiratory and skin disease risks from air pollution during the 2024 Beirut conflict}
{Disaster Med Public Health Prep}
{Satellite-derived PM$_{2.5}$ over Beirut's southern suburbs compared pre-war
(1 Aug--22 Sep 2024) against war (23 Sep--29 Nov 2024), plus ground PM$_{2.5}$ at five
sites over ten days. Satellite median rose \textbf{15.11 $\rightarrow$
20.21\,$\mu$g/m$^3$}; ground median in the direct war zone 19.33 $\rightarrow$
22.16\,$\mu$g/m$^3$. Applying BAPHE exposure--response functions gives \textbf{RR 1.011
(AR 1.13\%) respiratory} and \textbf{RR 1.032 (AR 3.10\%) skin} admissions, corroborated
by WHO AirQ+. The estimates are transferred, not observed --- no admissions data were
analysed --- so these are projections conditional on a borrowed concentration--response
function.}
{10.1017/dmp.2026.10412}

\paperentry{Li et al.\ 2026b}
{Burden trends of Alzheimer disease and other dementias in China, 1990--2023}
{Ann Acad Med Singap}
{GBD 2023 secondary analysis, GATHER-compliant. Deaths rose 122{,}411 $\rightarrow$
586{,}002; DALYs and prevalence more than tripled; age-standardised rates rose only
modestly, so the burden growth is demographic. Women carry higher ASRs in 2023 but male
rates are rising faster; burden concentrates at $\geq$75 years. Among the modelled risk
factors, \textbf{ambient particulate matter accounts for the largest attributable share
in 2023}, while household solid-fuel air pollution declined. That attributable fraction
is a GBD model output, not an epidemiological estimate from Chinese cohorts, and it
inherits the comparative risk assessment's exposure--response assumptions wholesale.}
{10.47102/annals-acadmedsg.2026146}

\paperentry{Vellaiyan 2026}
{Diesel NO$_x$ emissions and public health: a source--pathway--receptor synthesis}
{J Environ Health Sci Eng}
{Review spanning combustion science, atmospheric chemistry, exposure assessment and
epidemiology --- disciplines that normally publish past each other. Useful on near-road
gradients, street-canyon trapping, and the gap between certification cycles and
real-world driving emissions. A synthesis, not new evidence; the epidemiological half is
thinner than the combustion half.}
{10.1007/s40201-026-00995-z}

\paperentry{Makkar et al.\ 2026}
{Fungal biotechnology for air-pollution mitigation: mechanisms and engineered systems}
{Environ Sci Pollut Res}
{Proposes mycofiltration as a complement to conventional filtration on the argument that
HEPA-class media remove particulate mass efficiently but interact weakly with the
chemically reactive fraction (aldehydes, PAHs, VOCs, transition metals). Mechanisms
offered: extracellular oxidoreductases for aromatic VOC/PAH oxidation,
chitin--glucan--melanin wall polymers for metal and soot immobilisation, porous mycelium
for interception. Laboratory and pilot-scale evidence only --- no field deployment and no
bioaerosol-emission safety data. Speculative.}
{10.1007/s11356-026-38029-8}

\band{Neuro, cognition \& mental health\hfill 2 records}{Deep}

\paperentry{Calder\'on-Garcidue\~nas et al.\ 2026}
{Abnormal eye movements reflect early cortical and volumetric brain changes in PM$_{2.5}$-exposed urban youth}
{J Alzheimers Dis}
{Video-based eye tracking (EyeLink 1000-plus) in 80 volunteers aged 33$\pm$11 from
Metropolitan Mexico City and lower-exposure Cuernavaca, with structural MRI in 45 MMC
subjects. \textbf{Oculomotor variables did not differ significantly between cities} ---
the headline null. Abnormal antisaccades, low pursuit gain and square-wave jerks were
merely \emph{more frequent} in MMC, and the significant findings are correlations between
oculomotor measures and frontal--temporal--parietal cortical, hippocampal, thalamic,
caudate, amygdalar and cerebellar structure \emph{within} the exposed group. With
$n=45$ and that many regions the multiple-comparison exposure is severe, and no
correction is described. The group's long-running MMC series is valuable, but this
instalment does not establish an exposure contrast.}
{10.1177/13872877261471048}

\paperentry{Cattarinussi et al.\ 2026}
{Independent and joint effects of outdoor air pollution and genetic risk on adolescent mental health trajectories}
{medRxiv}
{ABCD Study, $n=10{,}620$, ages 9--16, sex-stratified trajectories of internalising
symptoms and psychotic-like experiences against residential PM$_{2.5}$ and polygenic
scores for schizophrenia and major depression. PM$_{2.5}$ was associated with
\textbf{smaller decreases in PLEs over time in females (p-FDR = 0.005)}, with no effect
on internalising trajectories in either sex, and a \textbf{PM$_{2.5}$ $\times$ PRS-MDD
interaction on PLE trajectories in females (p-FDR = 0.048)} --- a borderline interaction
that will not survive replication unless the effect is real. PM$_{2.5}$ is
residential-address-assigned, so within-person exposure misclassification is large and
non-differential, which biases toward the null for main effects but has unpredictable
sign for interactions. Preprint.}
{10.64898/2026.07.12.26357864}

\band{Cardiovascular \& metabolic\hfill 2 records}{Teal}

\paperentry{Luo et al.\ 2026}
{Short-term ambient air pollution and cardiovascular mortality: an 11-year time-series in Nanning}
{Front Public Health}
{87{,}913 CVD deaths, 2014--2024; generalized additive quasi-Poisson adjusted for
long-term trend, temperature, RH and day of week, across single-day lags 0--7 and
cumulative windows 0--1 to 0--7. At lag 0, \textbf{RR 1.014 (1.008--1.020) for
PM$_{2.5}$}, 1.009 (1.006--1.013) for PM$_{10}$ and 1.018 (1.011--1.026) for NO$_2$;
SO$_2$ peaks at lag 1 with \textbf{1.044 (1.018--1.069)}, the largest coefficient in the
set. Subtropical southwestern China is genuinely under-represented in this literature.
Two-pollutant models were run, but the paper does not resolve the SO$_2$/PM$_{2.5}$
collinearity its own results imply, so the apparent SO$_2$ dominance should not be read
as source-specific.}
{10.3389/fpubh.2026.1878398}

\paperentry{Chrysohoou et al.\ 2026}
{Preventing cardiovascular disease within the built and living environment: an expert opinion}
{Eur J Prev Cardiol}
{Expert opinion placing air and noise pollution alongside food environment, green space,
physical activity and psychosocial stress in a single built/living/social framework, with
inflammation, oxidative stress, neuroendocrine dysregulation and chronic stress as the
shared mechanistic substrate. No new data and no quantitative synthesis; its value is as
a citable framing document for grant and introduction text, nothing more.}
{10.1093/eurjpc/zwag387}

\band{Other clinical endpoints\hfill 2 records}{Slate}

\paperentry{Ouyang et al.\ 2026}
{Independent and joint effects of short-term air pollutant exposure on eczema in two Chinese cities}
{J Hazard Mater}
{879{,}356 daily eczema outpatient records from two Guangdong dermatology hospitals
matched to daily PM$_{2.5}$, PM$_{10}$, NO$_2$, SO$_2$ and O$_3$; distributed lag
non-linear models for single pollutants and quantile g-computation for two
four-pollutant mixtures. Per 10\,$\mu$g/m$^3$, \textbf{SO$_2$ gave the largest
single-pollutant effect at +22.63\%}. Each quartile increase in the PM$_{10}$-based
mixture raised visits 4.14\% (49.5\% of the weight on PM$_{10}$); the PM$_{2.5}$-based
mixture raised them 4.21\%, but with the weight on SO$_2$ (38.1\%) and NO$_2$ (37.5\%).
Both mixtures correspond to $\approx$18 extra visits per day. Outpatient visits confound
severity with care-seeking behaviour, and no confidence intervals are reported for the
mixture estimates.}
{10.1016/j.jhazmat.2026.143124}

\paperentry{Zurita et al.\ 2026}
{Childhood leukaemia and proximity to benzene emission sources in the Mexico City area}
{Pediatr Blood Cancer}
{Ecological study of acute lymphoblastic leukaemia cases from one paediatric hospital,
2010--2021, against counts of benzene-emitting establishments (petroleum and coal product
manufacture, fuel and lubricant retail, printing, chemical industry) within distance
buffers, using multivariate spatial models by municipality, sex and age group.
Significant positive association with source counts for ages 2--17. Single-hospital case
ascertainment plus municipality-level exposure is a textbook ecological design: no
individual exposure, no residential history, and referral patterns to one hospital are
almost certainly spatially structured. Hypothesis-generating at best --- the authors
say as much.}
{10.1002/1545-5017.70482}

\band{Respiratory \& allergic\hfill 1 record}{Sage}

\paperentry{Wang et al.\ 2026}
{Environmental and socioeconomic influences on childhood asthma and pneumonia burden}
{Front Med}
{Retrospective cohort from paediatric electronic medical records at one Wuhu hospital,
Anhui, January 2022--December 2025: 936 children aged 1--14 (47.4\% asthma, 36.5\%
pneumonia, 16.0\% both), with regional NO$_2$, PM$_{2.5}$, temperature and humidity and
census-based socioeconomic status. Higher NO$_2$ and PM$_{2.5}$ were associated with
asthma; biomass fuel use and poor household ventilation with pneumonia; lower SES
strongly with pneumonia but not asthma. Pneumonia peaked in winter, asthma in spring.
Hospital-based case series with area-level exposure and no comparison population --- the
associations describe who among the diagnosed had which condition, not incidence.}
{10.3389/fmed.2026.1872723}

\band{Reproductive \& developmental\hfill 1 record}{Sky}

\paperentry{Li et al.\ 2026}
{Ambient air pollutant exposure and semen quality: a refined multi-window panel study}
{Ecotoxicol Environ Saf}
{One-year prospective panel of 32 male undergraduates aged 18--22 with four quarterly
follow-ups spanning the 90-day spermatogenic cycle, using both daily-average and
cumulative-dose exposure metrics and --- the methodological contribution --- stratifying
the 15--69-day meiotic window into three biologically distinct sub-stages. Linear
mixed-effects, two-pollutant and Bayesian weighted quantile sum models. Per IQR of
full-cycle O$_3$, \textbf{sperm concentration fell 15.45\% (95\% CI $-$20.95 to
$-$9.58)}, with the sub-stage peak at days 31--45 ($-$13.30\%, $-$18.09 to $-$8.24);
NO$_2$ acted earliest (15--30 d), CO later (31--69 d). With $n=32$ the sub-stage
resolution outruns the power, and homogeneous undergraduates limit generalisability by
design. The window-stratification framework is worth borrowing regardless.}
{10.1016/j.ecoenv.2026.120581}

\band{Occupational \& indoor exposure\hfill 1 record}{Clay}

\paperentry{Gwak et al.\ 2026}
{Size-resolved modelling of airborne viral RNA copies in indoor environments}
{Risk Anal}
{Size-resolved framework over 1--2000\,$\mu$m coupling saliva-specific evaporation and
residue formation with gravitational settling, across temperature, RH and two ambient
pressures (101.325 and 75.3\,kPa), combined with saliva-mass uncertainty, viral load,
half-life, exhalation mode and ventilation. \textbf{Talking dominates airborne RNA
persistence} because coughing and sneezing generate larger droplets that settle quickly.
Predictions matched hospital air-sampling data within roughly one order of magnitude at
the 50th saliva-mass percentile. The framework's own caveat is the binding one: RNA decay
stands in for infectivity decay, and if viral RNA preferentially concentrates in
long-suspended fine aerosol the size-resolved conclusions shift. Directly relevant to
size-resolved optical sensing of bioaerosol in occupied spaces.}
{10.1111/risa.70318}

\band{Sensing, forecasting \& instrumentation\hfill 1 record}{Amber}

\paperentry{Aher \& Ravindran 2026}
{Comment on advancing air pollution forecasting: operational challenges beyond predictive accuracy}
{Environ Sci Pollut Res}
{Letter responding to a 2026 forecasting review. The substantive critique is that the
review optimises for predictive accuracy while under-treating what actually breaks
operational systems: model transferability into data-sparse regions, uncertainty
quantification, interpretability of deep architectures, and input data-quality
constraints. It points to adaptive mesh refinement, AI-assisted chemical transport models
and hybrid physics--ML frameworks as the productive direction. No new results --- but the
transferability-into-data-sparse-regions argument is precisely the problem low-cost
sensor networks exist to solve, and the letter states it cleanly.}
{10.1007/s11356-026-38087-y}

\clearpage

% ==================================================================== TABLE
\section*{\sffamily\bfseries\color{Deep}Corpus register}
\vspace{-1.5mm}{\color{Amber}\rule{\linewidth}{1.1pt}}\vspace{2.5mm}

{\fontsize{6.7}{8.4}\selectfont
\setlength{\tabcolsep}{2.4pt}
\begin{longtable}{@{}>{\RaggedRight}p{25mm}>{\RaggedRight}p{41mm}>{\RaggedRight}p{27mm}%
>{\RaggedRight}p{23mm}>{\RaggedRight}p{18mm}>{\RaggedRight}p{15mm}@{}}
\toprule
\textsf{\textbf{First author}} & \textsf{\textbf{Focus}} & \textsf{\textbf{Journal}} &
\textsf{\textbf{Design}} & \textsf{\textbf{Metric}} & \textsf{\textbf{Region}}\\
\midrule
\endfirsthead
\toprule
\textsf{\textbf{First author}} & \textsf{\textbf{Focus}} & \textsf{\textbf{Journal}} &
\textsf{\textbf{Design}} & \textsf{\textbf{Metric}} & \textsf{\textbf{Region}}\\
\midrule
\endhead
\multicolumn{6}{@{}l@{}}{\cellcolor{Mist}\textsf{\textbf{\textcolor{Deep}{Neuro / mental health}}}}\\[0.6mm]
\href{https://doi.org/10.1001/jamanetworkopen.2026.25015}{Dales} & Wildfire-sourced PM$_{2.5}$ and ED visits for migraine and other primary headache & \textit{JAMA Netw Open} & Case-crossover & PM$_{2.5}$ (wildfire) & Canada\\
\href{https://doi.org/10.1093/ageing/afag215}{Glover} & Umbrella meta-review of urban environmental exposures and cognitive health & \textit{Age Ageing} & Umbrella review / meta-analysis & PM$_{2.5}$ & Global\\
\href{https://doi.org/10.1016/j.nsa.2026.107019}{Goossens} & In utero air pollution exposure and pre-/postnatal brain development & \textit{Neurosci Appl} & Narrative review & PM$_{2.5}$ / PM$_{10}$ / PM$_{\mathrm{coarse}}$ & Global\\
\href{https://doi.org/10.1016/j.isci.2026.116617}{Ma} & PM$_{2.5}$-induced neurotoxicity and depression: exposure risks and mechanisms & \textit{iScience} & Narrative review & PM$_{2.5}$ + components & Global\\
\href{https://doi.org/10.3760/cma.j.cn112150-20250804-00752}{Pan} & Maternal PM exposure and offspring neurodevelopment: effects and mechanisms & \textit{Zhonghua Yu Fang Yi Xue Za Zhi} & Narrative review & PM (bulk) & China\\
\addlinespace[0.9mm]
\multicolumn{6}{@{}l@{}}{\cellcolor{Mist}\textsf{\textbf{\textcolor{Teal}{Cardiovascular \& metabolic}}}}\\[0.6mm]
\href{https://doi.org/10.3760/cma.j.cn112338-20251120-00843}{Huang} & PM$_{2.5}$ and dyslipidaemia in rural Chinese children; dietary effect modification & \textit{Zhonghua Liu Xing Bing Xue Za Zhi} & Cross-sectional (multistage) & PM$_{2.5}$ & China\\
\href{https://doi.org/10.32394/pe/218966}{Kazmierska} & Air pollution and cardiovascular disease: review of current data & \textit{Przegl Epidemiol} & Narrative review & PM$_{2.5}$ / PM$_{10}$ & Global\\
\href{https://doi.org/10.1007/s00125-026-06814-2}{Xi} & Air pollution, built environment and temperature vs. prediabetes phenotypes & \textit{Diabetologia} & Prospective cohort & PM$_{2.5}$ & Germany\\
\href{https://doi.org/10.1161/JAHA.125.046917}{Yang} & Short-term PM$_{2.5}$ and hospitalisation/ED use in atrial fibrillation & \textit{J Am Heart Assoc} & Case-crossover & PM$_{2.5}$ (wildfire-influenced) & USA\\
\addlinespace[0.9mm]
\multicolumn{6}{@{}l@{}}{\cellcolor{Mist}\textsf{\textbf{\textcolor{Sky}{Reproductive \& developmental}}}}\\[0.6mm]
\href{https://doi.org/10.1016/j.jes.2025.10.043}{Chen} & Prenatal air pollution, fetal growth trajectory phenotypes and child neurodevelopment & \textit{J Environ Sci (China)} & Prospective cohort & PM$_{2.5}$ / PM$_{10}$ & China\\
\href{https://doi.org/10.1016/j.ecoenv.2026.120566}{Huang} & Joint prenatal bisphenol + air-pollutant mixture exposure and preschool IQ & \textit{Ecotoxicol Environ Saf} & Prospective cohort (mixtures) & PM$_{2.5}$ / PM$_{10}$ & China\\
\href{https://doi.org/10.1007/s40572-026-00552-8}{Peterson} & Ambient air pollution, PAHs and pubertal development: critical review & \textit{Curr Environ Health Rep} & Critical review & PM$_{2.5}$ components & Global\\
\href{https://doi.org/10.1016/j.envres.2026.125313}{Xie} & Long-term PM$_{2.5}$ and timing of pubertal development in boys and girls & \textit{Environ Res} & Prospective cohort & PM$_{2.5}$ & China\\
\addlinespace[0.9mm]
\multicolumn{6}{@{}l@{}}{\cellcolor{Mist}\textsf{\textbf{\textcolor{Sage}{Respiratory \& allergic}}}}\\[0.6mm]
\href{https://doi.org/10.1007/s00484-026-03289-z}{Bayar Muluk} & Climate change and otolaryngology: airway disease impacts and emission reduction & \textit{Int J Biometeorol} & Narrative review & PM$_{2.5}$ & Global\\
\href{https://doi.org/10.4103/mgr.MEDGASRES-D-25-00164}{Liao} & Short-term air pollution, respiratory symptoms and lung function in asthma & \textit{Med Gas Res} & Panel study & PM$_{2.5}$ / PM$_{10}$ & China\\
\href{https://doi.org/10.3390/epidemiologia7040097}{Sanchez de Prada} & Ambient PM and asthma-coded admissions in children 0-3 y, Spain & \textit{Epidemiologia} & Ecological panel & PM$_{2.5}$ / PM$_{10}$ & Spain\\
\href{https://doi.org/10.1016/j.jhazmat.2026.143076}{Yang} & Allergenic vs. toxic risk of airborne seed hairs compared with ambient PM$_{2.5}$ & \textit{J Hazard Mater} & Field measurement & PM$_{2.5}$ + bioaerosol & China\\
\addlinespace[0.9mm]
\multicolumn{6}{@{}l@{}}{\cellcolor{Mist}\textsf{\textbf{\textcolor{Violet}{Mechanistic toxicology}}}}\\[0.6mm]
\href{https://doi.org/10.1016/j.jes.2025.12.074}{Ding} & PM$_{2.5}$ promotes vascular calcification via miR-27b-5p/HES1 signalling & \textit{J Environ Sci (China)} & In vitro + organoid + murine & PM$_{2.5}$ & China\\
\href{https://doi.org/10.1007/s00360-026-01698-5}{Maraschi} & Osmoregulatory disruption and metal bioaccumulation from settleable atmospheric PM & \textit{J Comp Physiol B} & Controlled exposure (ecotox) & Settleable APM & Brazil\\
\href{https://doi.org/10.1016/j.jes.2025.11.018}{Wu} & Macrophage NLRP3 inflammasome in real-ambient PM-induced hepatic glucose dysregulation & \textit{J Environ Sci (China)} & Real-ambient murine + co-culture & Real-ambient PM & China\\
\addlinespace[0.9mm]
\multicolumn{6}{@{}l@{}}{\cellcolor{Mist}\textsf{\textbf{\textcolor{Clay}{Occupational \& indoor}}}}\\[0.6mm]
\href{https://doi.org/10.3389/fmedt.2026.1796828}{Chen} & Nested suction instrument to cut intraoperative surgical-smoke PM$_{2.5}$ & \textit{Front Med Technol} & Bench / device experiment & PM$_{2.5}$ (surgical smoke) & China\\
\href{https://doi.org/10.1093/annweh/wxag055}{Cole} & Determinants of respirable crystalline silica exposure in Australian tunnelling & \textit{Ann Work Expo Health} & Observational exposure study & Respirable dust / RCS & Australia\\
\href{https://doi.org/10.1016/j.jhazmat.2026.143081}{Ding} & Stage-resolved hair Zn-isotope dynamics as a firefighting combustion-exposure tracer & \textit{J Hazard Mater} & Repeated-measures biomarker & Combustion-derived PM & China\\
\href{https://doi.org/10.3390/jox16040119}{Ramsden} & Xenobiotic hazards in aircraft cabin air & \textit{J Xenobiot} & Review + risk economics & Ultrafine particles & Global\\
\addlinespace[0.9mm]
\multicolumn{6}{@{}l@{}}{\cellcolor{Mist}\textsf{\textbf{\textcolor{Amber}{Sensing, forecasting \& instrumentation}}}}\\[0.6mm]
\href{https://doi.org/10.1016/j.jenvman.2026.130515}{Shi} & Adaptive spatiotemporal graph transformer for multi-site multi-step PM$_{2.5}$ forecasting & \textit{J Environ Manage} & Deep-learning model & PM$_{2.5}$ & China / India\\
\href{https://doi.org/10.3390/bios16070380}{Tsuchiya} & Seasonal and indoor-outdoor characteristics of fluorescent aerosol particles in offices & \textit{Biosensors} & Real-time sensor campaign & Bioaerosol / coarse & Japan\\
\addlinespace[0.9mm]
\multicolumn{6}{@{}l@{}}{\cellcolor{Mist}\textsf{\textbf{\textcolor{Amber}{Exposure assessment \& modelling}}}}\\[0.6mm]
\href{https://doi.org/10.1029/2025GH001621}{Solorzano-Araque} & Multiscale exposure modelling: LOTOS-EUROS + mobility surveys + particle toxicity & \textit{GeoHealth} & CTM + mobility integration & PM$_{2.5}$ / PM$_{10}$ & Colombia\\
\addlinespace[0.9mm]
\multicolumn{6}{@{}l@{}}{\cellcolor{Mist}\textsf{\textbf{\textcolor{Coral}{Burden, policy \& mitigation}}}}\\[0.6mm]
\href{https://doi.org/10.1002/hsr2.72886}{Karamzad} & Ischaemic heart disease burden in MENA 1990-2021; PM attributable fraction 28.2\% & \textit{Health Sci Rep} & GBD secondary analysis & Ambient PM & MENA\\
\href{https://doi.org/10.1371/journal.pone.0347989}{Rees} & HEPA filtration to reduce respiratory infection in care homes & \textit{PLoS One} & RCT process evaluation & Filtered PM & UK\\
\href{https://doi.org/10.1016/j.jes.2026.01.009}{Shan} & Policy-driven emission reductions and post-pandemic rebound, high-resolution inventories & \textit{J Environ Sci (China)} & Emission inventory & PM emissions & China\\
\addlinespace[0.9mm]
\multicolumn{6}{@{}l@{}}{\cellcolor{Mist}\textsf{\textbf{\textcolor{Slate}{Other clinical endpoints}}}}\\[0.6mm]
\href{https://doi.org/10.3390/dj14070440}{Lupita} & PM$_{2.5}$, preventive-care capacity and caries/edentulism burden across the EU & \textit{Dent J} & Ecological panel & PM$_{2.5}$ & EU-27\\
\href{https://doi.org/10.1016/j.jdin.2026.06.008}{Park} & Long-term PM exposure and risk of androgenetic alopecia & \textit{JAAD Int} & Nationwide cohort & PM$_{2.5}$ / PM$_{10}$ & South Korea\\
\href{https://doi.org/10.3171/2026.3.PEDS2615}{Sun} & Environmental exposure and area deprivation vs. paediatric hydrocephalus outcomes & \textit{J Neurosurg Pediatr} & Registry cohort + geospatial & PM$_{2.5}$ (satellite-derived) & USA\\
\addlinespace[0.9mm]
\bottomrule
\end{longtable}}


\vspace{4mm}
\begin{keybox}
{\sffamily\bfseries\footnotesize\color{Deep}Provenance and method}\\[1.2mm]
{\fontsize{7.2}{8.8}\selectfont
Entry window \textbf{30 July 2026}, contiguous with the 29 July issue --- no gap.
Sources: \textbf{PubMed} (E-utilities and connector) queried on the \texttt{[Date -
Entry]} field along two independent axes, health (16 hits) and instrumentation (66 hits);
\textbf{Europe PMC} \texttt{CREATION\_DATE:[2026-07-30 TO 2026-07-30]} (5 hits, 3 of them
duplicates of PubMed records, 2 new preprints); \textbf{arXiv} relevance query with a
local recency screen (60 entries, none newer than 29 July, none in scope);
\textbf{OpenAlex} returned HTTP 429 and remains unusable for date-windowed harvesting;
\textbf{ClinicalTrials.gov} returned zero registrations or updates in the window.
\textbf{18 records in scope.}
\\[1.0mm]
\textbf{Two rejection groups deserve explicit statement.} (i) The 66-hit instrumentation
axis was almost entirely recall noise --- \texttt{calibration}, \texttt{machine learning}
and \texttt{monitoring} match clinical prediction-model papers --- and yielded two
in-scope records. (ii) The \textbf{Consensus} sweep, which exists precisely because
PubMed does not index \textit{Atmos.\ Meas.\ Tech.} or \textit{Atmos.\ Chem.\ Phys.},
\emph{has no entry-date filter}. Its eight low-cost-sensor hits resolved via Crossref to
publication dates between 20 January and 6 July 2026, and one
(\texttt{10.1080/02786826.2026.2676293}) was already carried in the 29 July issue. All
eight were rejected to \texttt{state/rejected.jsonl} rather than presented as new. The
structural consequence is that \textbf{this pipeline currently has no date-windowed
channel into the atmospheric-measurement literature}; closing it needs a Crossref
by-journal feed for AMT/ACP keyed on \texttt{created}, not another Consensus query.
\\[1.0mm]
Subtopic, design, particle-metric and geography labels were assigned from title,
abstract, keywords and MeSH terms; assignment is single-label by dominant contribution
and is therefore a judgement, not a controlled vocabulary. Effect estimates in the forest
plot are transcribed verbatim from abstracts and use heterogeneous exposure contrasts ---
see the figure caption. Two records are Research Square preprints and one is a medRxiv
preprint; none has been peer reviewed, and each is marked as such in the digest.
All DOI links resolve to the publisher of record. Content derived from PubMed and
Europe PMC metadata; abstracts remain the copyright of their respective publishers.\par}
\end{keybox}

\end{document}
